% -------------------------------------------------------------------------
% WBS ID: RES-GEO-SPEC
% Deliverable: Integrated Barrier-Class Geometry Specification
% Status: In progress
% Evidence status: Active barrier-class comparison scope established
% Review status: Not reviewed
% Last updated: 2026-07-19
% Dependencies: RES-GEO-TAX through RES-GEO-CST
% Downstream interfaces: RES-MOD-IBC, RES-NUM-MSH, RES-PHY-BAR, RES-PHY-LOD, RES-VER-MET, Software geometry manifest
% -------------------------------------------------------------------------

\section{RES-GEO-SPEC --- Integrated Barrier-Class Geometry Specification}
\label{sec:res-geo-spec}

\subsection{Purpose and Scope}
\label{sec:res-geo-spec-purpose}

This Deliverable consolidates the active fixed-rigid barrier geometry
specification. It preserves the WBS ID while changing the active output
away from geometry generation and toward reproducible wall-type and
obstacle/dissipating/array comparisons in the two Tohoku corridors.

\subsection{Research Questions}
\label{sec:res-geo-spec-research-questions}

\begin{enumerate}
  \item Are the active geometry classes, descriptors, placement rules
    and representations sufficient for local 3D hydrodynamic
    comparison?
  \item Which fields must the Software workstream implement in the
    geometry manifest before data handling can proceed?
  \item Which geometry ideas are explicitly deferred?
\end{enumerate}

\subsection{Terminology and Definitions}
\label{sec:res-geo-spec-definitions}

% TODO: Define the technical terminology, symbols, quantities and
% classifications used in this Deliverable.

\subsection{Literature Evidence}
\label{sec:res-geo-spec-literature-evidence}

\textcite{deljesus2012porous} supports the impermeable/opening-aware
barrier distinction and limits unresolved porous claims.
\textcite{watanabe2022validation} supports preserving geometry and
output identifiers for local impact validation. \textcite{mori2012nationwide}
supports retaining the two Tohoku corridor settings as contrasting
hydrodynamic contexts. Project conclusion: geometry specification is a
data and modelling contract for comparison, not a literature survey or
automatic design generator.

\subsection{Governing Relations and Quantitative Definitions}
\label{sec:res-geo-spec-governing-relations}

The minimum geometry-case contract is:

\begin{align}
  \mathcal{C}^{\mathrm{geo}}_j
  &=
  \left\{
    \mathtt{corridor\_id},\,
    \mathtt{geometry\_id},\,
    \mathcal{G}_j,\,
    \boldsymbol{\theta}^{\mathrm{geo}}_j,\,
    \Gamma_{\mathrm{B},j},\,
    \mathcal{P}_j
  \right\}
  \label{eq:res-geo-spec-case-contract}
\end{align}

where $\mathcal{G}_j$ is the descriptor record,
$\boldsymbol{\theta}^{\mathrm{geo}}_j$ is the parameter vector,
$\Gamma_{\mathrm{B},j}$ is the fixed barrier patch and
$\mathcal{P}_j$ records provenance, CRS/datum linkage and unresolved
physics flags. Source role: this equation integrates the decisions in
RES-GEO-TAX, RES-GEO-PAR and RES-GEO-REP. Limitation: it does not
define solver implementation or validation results.

\subsection{Geometry Family or Representation}
\label{sec:res-geo-spec-geometry-family-representation}

The active representation is fixed terrain plus fixed rigid barrier
patches. Wall-type and obstacle/array-type cases share the same
manifest structure so they can be compared through common
hydrodynamic metrics.

\subsection{Design Variables and Parameter Bounds}
\label{sec:res-geo-spec-design-variables-parameter-bounds}

Only descriptor-level comparison variables are active. Generative
control points, topology optimisation variables and learned geometry
latents are deferred.

\subsection{Geometric Descriptors}
\label{sec:res-geo-spec-geometric-descriptors}

Minimum descriptors are class, position, orientation, crest elevation,
height, width/thickness, footprint length, spacing, row count, opening
fraction, roughness flag, blockage ratio and provenance quality.

\subsection{Placement and Arrangement Rules}
\label{sec:res-geo-spec-placement-arrangement-rules}

Placement shall be corridor-relative and shall preserve fixed
measurement sections across class comparisons.

\subsection{Feasibility Constraints}
\label{sec:res-geo-spec-feasibility-constraints}

Feasibility requires terrain coverage, meshability, stable wall patch
labels and no unresolved physics dependency for the active baseline.

\subsection{Two- and Three-Dimensional Representation}
\label{sec:res-geo-spec-two-three-dimensional-representation}

The local three-dimensional representation is a fixed barrier surface
used for no-slip/no-penetration boundary conditions and for pressure,
shear, force, impulse and moment extraction.

\subsection{Candidate Approaches}
\label{sec:res-geo-spec-candidate-approaches}

% TODO: Define the credible candidate approaches considered.

\subsection{Critical Comparison}
\label{sec:res-geo-spec-critical-comparison}

% TODO: Compare candidates using physically and project-relevant criteria.
% Do not select an approach solely on popularity or implementation ease.

\subsection{Assumptions and Applicability}
\label{sec:res-geo-spec-assumptions}

% TODO: State assumptions and the conditions under which they are valid.

\subsection{Limitations and Failure Conditions}
\label{sec:res-geo-spec-limitations}

% TODO: State where the evidence, model or selected approach ceases to be
% reliable or applicable.

\subsection{Project-Specific Conclusions}
\label{sec:res-geo-spec-conclusions}

The integrated RES-GEO output now supports the immediate project
scope: Tohoku/local comparison and barrier-class hydrodynamic
comparison. It does not authorise geometry generation, structural
damage simulation or ML surrogate training.

\subsection{Selected Approach and Rationale}
\label{sec:res-geo-spec-selected-approach}

The active approach is descriptor-first fixed-rigid barrier-class
comparison. Existing-defence reconstruction may be added where source
geometry is available. Unrestricted generative geometry, topology
optimisation, surrogate-led geometry search and deformable/damage
geometry are deferred until hydrodynamic extraction and validation
are stable.

\subsection{Unresolved Decisions}
\label{sec:res-geo-spec-unresolved-decisions}

\begin{itemize}
  \item TODO: citable Kamaishi bay-mouth breakwater geometry and 2011
    condition source.
  \item Final corridor-relative placement sections and geometry bounds.
  \item Geometry file format and HDF5/structured schema details.
  \item Roughness/opening/porous-closure activation policy.
\end{itemize}

\subsection{Requirements and Handoffs}
\label{sec:res-geo-spec-requirements}

Software may begin the geometry-manifest MVP in parallel with
Research: parse geometry case records, generate corridor-relative
barrier patches, preserve patch identifiers, run clipping/meshability
checks and export structured geometry metadata for replay-boundary
and local 3D scaffolds. This is a data-handling handoff only, not
solver implementation.

\subsection{Deliverable Completion Record}
\label{sec:res-geo-spec-completion-record}

% TODO: Record whether the Deliverable is:
%
% - Not started
% - In progress
% - Draft complete
% - Reviewed
% - Accepted
%
% Acceptance requires:
%
% - the research questions to be answered;
% - claims to be supported by appropriate evidence;
% - assumptions and limitations to be explicit;
% - the project-specific conclusion to be stated;
% - downstream requirements to be traceable;
% - unresolved decisions to be closed or formally deferred.
